\begingroup
\definecolor{zonCodeBackground}{HTML}{FAF5F0}
\definecolor{zonCodeRule}{HTML}{B08968}
\definecolor{zonCodeKeyword}{HTML}{B64080}
\definecolor{zonCodeComment}{HTML}{008000}
\begin{figure}[t]
\refstepcounter{lstlisting}
\label{fig:zon_mlp_listing}
\begin{center}
Listing~\thelstlisting: Minimal PyTorch implementation of the ZoN MLP training forward pass.
\end{center}
\begin{tcblisting}{
    enhanced,
    listing only,
    colback=zonCodeBackground,
    colframe=zonCodeRule,
    boxrule=0.6pt,
    arc=8pt,
    boxsep=0pt,
    left=8pt,right=8pt,top=8pt,bottom=8pt,
    listing options={
        language=Python,
        basicstyle=\fontencoding{T1}\fontfamily{pcr}\selectfont\small,
        keywordstyle=\color{zonCodeKeyword}\bfseries,
        commentstyle=\color{zonCodeComment},
        stringstyle=\color{purple},
        frame=none,
        aboveskip=0pt,belowskip=0pt,
        showstringspaces=false,
        columns=fullflexible,
        keepspaces=true,
        upquote=true,
        numbers=none
    }
}
def mlp_zon(x, W_k, W_v, act_fn, a, beta=1.0, l=-0.01):
    """
    x: [batch * seq_len, d_model]
    K: [d_model, d_inter]
    V: [d_inter, d_model]
    a: gate network, d_model -> d_inter
    """
    log_alpha = a(x) # shape: [batch * seq_len, d_inter]

    # Sample gates
    u = torch.empty_like(log_alpha).uniform_(eps, 1 - eps)
    s = torch.sigmoid((torch.logit(u) + log_alpha) / beta)
    # stretch beyond zero and clamp to 0.0
    z = (s + l * (1 - s)).clamp(min=0.0)
    # Expected number of active gates per token
    num_active = torch.sigmoid(log_alpha - beta * math.log(-l)).sum(-1)

    h = act_fn((x @ W_k) * z)
    o = h @ W_v

    return o, num_active
\end{tcblisting}
\end{figure}
\endgroup
